% =============================================================================
% 6 Evaluation. Written 2026-10-21 (Ben: "let's continue to evaluation"). The
% paper is published before any GH200 run (Ben, 2026-10-19), so s6 shows that
% the capacity model is accurate on A100 evidence and states what it projects,
% with its sensitivity. Numbers: paper/model/validate_a100.py (A100 checks) and
% paper/model/gain_curve.py point(..., overrides) (projection); analysis in
% paper_versions/m22_analysis/M22_ANALYSIS.md.
% =============================================================================

\section{Evaluation}\label{sec:eval}

\sysname{} is meant for a GH200; the measurements in this paper are from an
A100. The evaluation therefore asks three questions. Does the capacity model
predict the measured step time? Does its decision admit the right number of
riders? And what does the model, so
tested, project for the target, and how much does that projection depend on
each of its inputs?

\subsection{Setup}

The rig is that of \S\ref{sec:ablation}: one A100-80 (PCIe, 25.1\,GB/s measured
host to device) serving Qwen2.5-32B in fp16 as a decode node, with the
handover emulated and every step checked to be pure decode. The model's
profile comes from runs with the transfer's bytes removed: a single resident,
contexts 512 to 14\,724, up to five riders. It gives
$c_{\mathrm{m}} = 0.50$\,ms, the attention slope of \cref{eq:cS}, and
$F = 5.0$\,ms. The contention coefficient comes from the microbenchmark of
\cref{fig:abl-contention}; on A100 the copies arrive over PCIe, so the model
uses the PCIe value, $\kappa = 4.55$.

% Section 6.2: the model's prediction of the step overhead for runs it was not
% fitted on. Points: real-bytes sweeps at ctx 1024, 1 resident (M22c job
% 2190742, six runs; M23 job 2190851, two runs), and the matching bytes-removed
% sweeps (M23, two runs); figs/data/eval_dt.csv, rows 0-31 real, 32-39 removed.
% Lines: the model, with the profile from the bytes-removed ladder (M22, job
% 2189932: c_m 0.497, 0.2121 ms per 1k tok, F 5.0) and kappa from the
% microbenchmark at the PCIe rate (4.55). paper/model/validate_a100.py.
\begin{figure}[t]
\centering
\begin{tikzpicture}
\begin{axis}[
  width=0.95\columnwidth, height=4.6cm,
  xmin=0.6, xmax=4.4, ymin=4, ymax=12,
  xtick={1,2,3,4},
  xlabel={riders $r$},
  ylabel={step overhead $\Delta T$ (ms)},
  tick label style={font=\sffamily\scriptsize},
  label style={font=\sffamily\scriptsize},
  axis line style={black!60}, tick style={black!60},
  ymajorgrids, grid style={black!10},
  legend style={font=\sffamily\scriptsize, draw=none, fill=white,
                fill opacity=0.85, text opacity=1, at={(0.5,-0.30)},
                anchor=north, row sep=-1pt, legend columns=2,
                column sep=6pt},
  legend cell align=left,
  every axis plot/.append style={line width=0.8pt},
]
\addplot[ddrD, domain=1:4, samples=2] {5.0+1.512*x};
\addlegendentry{model, real bytes}
\addplot[only marks, ddrD, mark=*, mark size=1.3pt, opacity=0.7]
  table[x=r, y=dT, col sep=comma,
        restrict expr to domain={\coordindex}{0:31}] {figs/data/eval_dt.csv};
\addlegendentry{measured (8 runs)}
\addplot[hbmD, domain=1:4, samples=2] {5.0+0.714*x};
\addlegendentry{model, no transfer terms}
\addplot[only marks, hbmD, mark=o, mark size=1.6pt]
  table[x=r, y=dT, col sep=comma,
        restrict expr to domain={\coordindex}{32:39}] {figs/data/eval_dt.csv};
\addlegendentry{bytes removed (2 runs)}
\end{axis}
\end{tikzpicture}
\caption{\textbf{The model predicts runs it was not fitted on.} A100-80,
Qwen2.5-32B, context 1024, one resident, link loaded up to 98\% of the step.
The profile comes from a different sweep (bytes removed, contexts 512--14\,724)
and $\kappa$ from the microbenchmark of \cref{fig:abl-contention}. With real
bytes the prediction is within 1.7\% of the measured step; the same profile
without the contention and first-layer terms falls on the bytes-removed line,
up to 6\% low.}
\label{fig:eval-pred}
\end{figure}


\subsection{Does the Model Predict the Step?}\label{sec:eval-pred}

We compare the prediction with eight runs with real bytes at context 1024,
one resident and one to four riders, from two jobs on different days: 32 step
times with the link carrying up to 98\% of a step. \Cref{fig:eval-pred} shows them against
the model's prediction. The mean error is $-0.5\%$ and the largest is 1.7\%
(0.9\,ms); the throughput gain each point implies is predicted within 1.7\%.
Both transfer terms are needed. Without the contention and first-layer terms
of \cref{eq:cS} the same profile predicts the bytes-removed runs, which it
matches within 1.1\%, and falls up to 6.1\% below the real one.

At sixteen residents and contexts 128 to 512 the model predicts the per-rider
cost $c$ with real bytes within 0.14\,ms (0.49, 0.47 and 0.83\,ms predicted
against 0.41, 0.56 and 0.97 measured), starting from that configuration's own
bytes-removed runs, whose $c$ itself varies by up to 0.3\,ms between repeats.
The baseline step $T_0$ computed from \cref{eq:cS} is within 0.6\% of the
measured one with one resident, at both 1024 and 14\,724 tokens, but 5\% low
with sixteen: per-request work beyond the weight and KV reads is not in the
equation. This is why the method measures $T_0$ on the node rather than
computing it (\S\ref{sec:opoint}), and why the projection below varies it.
% src: validate_a100.py: real n=32 err mean -0.48% max 1.65% (0.92 ms); gain
%      err max 1.68%; bytes removed n=8 max 1.05%; without the two terms mean
%      -4.11% max 6.12%. 16 residents (M20, paper_runs/20260930/overhead/
%      22_c_context): MB c 0.394/0.269/0.430, predicted 0.493/0.469/0.829,
%      measured CC 0.413/0.559/0.967 (ctx 128/256/512); MB_c128 repeats 0.24
%      and 0.54. T0: predicted 45.37 vs 45.20 (B1 c1024), 48.27 vs 47.96-48.02
%      (B1 c14724), 45.59 vs 47.84 (B16 c128), 46.89 vs 49.50 (B16 c512).
%      kappa_PCIe = 0.235 %/(GB/s) x 1935 GB/s / 100 = 4.55 (E1 h2d).

\subsection{Does the Decision Admit the Right Riders?}\label{sec:eval-knee}

The decision of \cref{eq:decision} admits riders up to the knee,
$\lfloor T_0\Xlink/(S\btok)\rfloor$. At context 1024 with one resident the
knee is 4.09, so it admits four. Past it, a run to six riders shows where
the transfer stops hiding: four riders lie on the line of \cref{eq:law}, the
fifth adds 4\,ms above it, and from there the step is the link time plus
7\,ms. The measured knee is therefore at about 4.6 riders, and the rule's
choice is conservative by 12\%: it never admits a rider whose transfer is
exposed. That caution has a price here. Throughput at four riders is 8\% below
the best measured count, five; six riders lower it again.
% src: M22c (job 2190742, 6 runs): F 5.11-5.21, c 1.53-1.56 over r=1..4;
%      excess over the line at r=5 +3.8..+4.2 ms; step - link 6.8 (r=5),
%      7.5-7.6 (r=6); knee from link*r + 6.8 = T0 + F + c*r -> 4.58; decode
%      tok/s 88.4-88.8 (r=4), 96.6-96.7 (r=5), 94.7-94.8 (r=6).

\subsection{What the Model Projects}\label{sec:eval-proj}

\begin{table}[t]
\centering\small\setlength{\tabcolsep}{4pt}
\caption{\textbf{Sensitivity of the projection} at the operating point of
\S\ref{sec:opoint} (Qwen2.5-72B, fp8 weights, 16-bit KV, GH200, context
14\,724, $B{=}3$, $r{=}2$): each row changes one input of the model.}
\label{tab:sens}
\begin{tabular}{@{}l r r r@{}}
\toprule
Input changed & $c$ (ms) & gain & $\eta$ \\
\midrule
none (the model as stated)                  & 3.93 & $1.20\times$ & 1.9 \\
copies cost HBM nothing, $\kappa = 0$         & 2.65 & $1.27\times$ & 2.6 \\
$\kappa = 0.89$ (DAK's GH200 figure)         & 3.73 & $1.21\times$ & 2.0 \\
$\kappa = 1.5$                                & 4.46 & $1.17\times$ & 1.6 \\
weight-read time $\pm 5\%$                    & 3.93 & $1.18$--$1.21\times$ & 1.8--2.0 \\
KV read at 50\% of peak, not 64\%             & 4.46 & $1.19\times$ & 1.8 \\
$F = 3.6$\,ms (16-resident profile)           & 3.93 & $1.22\times$ & 2.1 \\
$F = 6.3$\,ms (1-resident profile)            & 3.93 & $1.15\times$ & 1.4 \\
\bottomrule
\end{tabular}
\end{table}
% src: gain_curve.py point(14724, ...): base c 3.932 G 1.199 eta 1.90;
%      kappa 0: 2.649/1.274/2.62 (first layer kept); 0.89: 3.730/1.210/2.01
%      (dak s2.3: up to 10% at 450/4000 = 11.25% -> 0.89); 1.5: 4.457/1.170/
%      1.62; t0_scale 0.95/1.05: 1.184/1.212, eta 1.76/2.02; eff_kv 0.50:
%      4.457/1.187/1.79; F 2.9*80/64 = 3.63: 1.216/2.06; F 5.0*80/64 = 6.25:
%      1.145/1.39. No row changes B (3) or r (2).

Applied to the target with the GH200's bandwidths, the model gives
$1.20\times$ at context 14\,724, and $\eta = 1.9$: the link's 10.5\% of added
bandwidth buys 20\% more throughput (\S\ref{sec:opoint}).
\Cref{tab:sens} varies each input over the range our measurements leave open.
The gain stays between $1.15\times$ and $1.27\times$, and $\eta$ between 1.4
and 2.6, above the bound of 1 that using the link to read the same requests
faster cannot exceed. No row changes the number of residents or riders. The
input that matters most is the fixed cost: $F$ measured with sixteen residents
is 40\% lower than with one, and the target's three residents lie between.
The contention coefficient comes next; DAK's figure for GH200 background
copies~\cite{dak} lands within 0.01 of the stated gain.

% Section 6.4: projected gain of each design class at the operating point
% (Qwen2.5-72B fp8, 16-bit KV, GH200, ctx 14724, B=3; r chosen by
% eq. decision for each design, 2026-10-22). Ben, 2026-10-21:
% "projected graph comparison with DirectKV/Pie/DAK". Values printed by
% paper/model/compare_designs.py; each is the capacity model applied to the
% mechanism the paper describes, not a model of the system.
\begin{figure}[t]
\centering
\begin{tikzpicture}
\begin{axis}[
  width=0.72\columnwidth, height=4.2cm,
  xbar, bar width=5pt,
  xmin=0.8, xmax=2.0,
  ymin=0.4, ymax=4.6,
  ytick={1,2,3,4},
  yticklabels={HBM only, aggregation (DAK), {zero-copy (DirectKV)},
               \sysname{}},
  xtick={0.8,1.0,1.2,1.4,1.6,1.8},
  xticklabel style={/pgf/number format/.cd,fixed,fixed zerofill,precision=1},
  xlabel={projected throughput gain over HBM only},
  tick label style={font=\sffamily\scriptsize},
  label style={font=\sffamily\scriptsize},
  axis line style={black!60}, tick style={black!60},
  xmajorgrids, grid style={black!10},
  legend style={font=\sffamily\scriptsize, draw=none, fill=white,
                fill opacity=0.85, text opacity=1, at={(0.5,-0.32)},
                anchor=north, legend columns=2, column sep=6pt},
  legend cell align=left,
  nodes near coords={\pgfmathprintnumber[fixed,fixed zerofill,precision=2]{\pgfplotspointmeta}},
  nodes near coords style={font=\sffamily\tiny, black!80},
  point meta=x,
]
\addplot[draw=hbmD, fill=hbmD!35, bar shift=3pt]
  coordinates {(1.000,1) (1.105,2) (0.897,3) (1.499,4)};
\addlegendentry{mechanism only}
\addplot[draw=ddrD, fill=ddrD, bar shift=-3pt]
  coordinates {(1.368,4)};
\addlegendentry{with \sysname{}'s measured overheads}
\draw[black!60, densely dashed, line width=0.6pt]
  (axis cs:1.794,0.4) -- (axis cs:1.794,4.6);
\node[font=\sffamily\tiny, black!70, anchor=south west, align=left]
  at (axis cs:1.81,1.6) {link\\ceiling,\\no costs\\(1.79)};
\end{axis}
\end{tikzpicture}
\caption{\textbf{Projected gain by design}, at the operating point of
\S\ref{sec:opoint} on a GH200 (context 14\,724, three residents). Each bar is
the capacity model applied to the mechanism each paper describes, not a model
of that system. Aggregation reads the same batch faster and is bounded by
$1+\Xlink/BW$. Zero-copy of every request's KV is bounded by the link, $\Xlink/(S\btok)$
requests per unit time. \sysname{} adds requests over the link, with the
rider count of \cref{eq:decision}. The light bar omits implementation
overheads; the dark bar adds \sysname{}'s measured $c_{\mathrm{m}}$ and
$F$.}
\label{fig:compare}
\end{figure}


\Cref{fig:compare} places this projection beside the other ways of using the
link, each modelled from the mechanism its paper describes. Reading the same
batch over HBM and the link together~\cite{dak} is bounded by
$1+\Xlink/BW = 1.105\times$, even if its copies cost nothing. Keeping every
request's KV in host memory and reading it over the link~\cite{directkv}
serves at most $\Xlink/(S\btok)$ requests per unit time, 0.90 of what HBM alone
serves at this context. Swapping layers~\cite{pie} grows the batch under the
same per-step ceiling as riders, $T_0\Xlink$, but swaps each layer back out
after use, so every moved byte crosses HBM twice; that costs it 0.06 to 0.08
of gain against \sysname{}. The ceiling itself, $1.67\times$, is what any design
that adds KV over the link could reach if moving and attending it cost nothing.
% src: paper/model/compare_designs.py: aggregation 1.105, zero-copy 0.897,
%      swap 1.295 / 1.139 (with c_m, F), HeteroKV 1.373 / 1.199, ceiling 1.667.
%      dak footnote 1 (bandwidth = HBM + min(C2C, DRAM)); directkv s6 l.1030
%      ("the manager stores the initial K,V entries from each layer in CPU
%      memory"); pie s3.1 (t_swap*m <= t_compute*n; "recently completed layers
%      ... can be swapped out to CPU memory").

\subsection{Under Link Contention}\label{sec:eval-pie}

% Section 6.5: gain under background link traffic (Ben, 2026-10-22: "put the
% graph show the potential degradation of Pie due to his write back"; the
% host-to-device control is shown with it). Medians over the 36 steady-state
% workloads of paper/model/steady_bg.py -> model/out/steady_bg.jsonl
% (overheads charged to both). d2h: het 1.180 flat; pie 1.136 1.114 1.085
% 1.031. h2d: het 1.180 1.156 1.103 1.025; pie 1.136 1.136 1.102 1.044.
\begin{figure}[t]
\centering
\begin{tikzpicture}
\begin{axis}[
  width=0.95\columnwidth, height=4.4cm,
  xmin=0, xmax=0.5, ymin=1.0, ymax=1.2,
  xtick={0,0.1,0.25,0.5},
  xticklabels={0, 10\%, 25\%, 50\%},
  ytick={1.0,1.05,1.1,1.15,1.2},
  yticklabel style={/pgf/number format/.cd,fixed,fixed zerofill,precision=2},
  xlabel={share of one link direction taken by other traffic},
  ylabel={gain over HBM only (median)},
  tick label style={font=\sffamily\scriptsize},
  label style={font=\sffamily\scriptsize},
  axis line style={black!60}, tick style={black!60},
  ymajorgrids, grid style={black!10},
  legend style={font=\sffamily\scriptsize, draw=none, fill=none,
                at={(0.5,-0.36)}, anchor=north, legend columns=2,
                column sep=8pt},
  legend cell align=left,
  mark size=1.6pt,
]
\addplot[ddrD, thick, mark=*]
  coordinates {(0,1.180) (0.1,1.180) (0.25,1.180) (0.5,1.180)};
\addlegendentry{\sysname{}, device to host}
\addplot[ddrD, thick, dashed, mark=o, mark options={solid}]
  coordinates {(0,1.180) (0.1,1.156) (0.25,1.103) (0.5,1.025)};
\addlegendentry{\sysname{}, host to device}
\addplot[black!60, thick, mark=square*]
  coordinates {(0,1.136) (0.1,1.114) (0.25,1.085) (0.5,1.031)};
\addlegendentry{layer swap (Pie), device to host}
\addplot[black!60, thick, dashed, mark=square, mark options={solid}]
  coordinates {(0,1.136) (0.1,1.136) (0.25,1.102) (0.5,1.044)};
\addlegendentry{layer swap (Pie), host to device}
\end{axis}
\end{tikzpicture}
\caption{\textbf{Gain under background link traffic.} Simulated decode on a
GH200 in steady state, median over 36 workloads (Qwen2.5-72B with 16-bit or
fp8 KV and Qwen2.5-32B, four context distributions, outputs of 128 to 2048
tokens).
Other work takes a share of one link direction. Layer swap writes every
swapped layer back, so it needs both directions and loses gain to traffic in
either. Riders read host to device and write back only the new token, so
traffic in the other direction leaves them unchanged, while traffic in their
own direction costs them more than it costs layer swap.}
\label{fig:linkload}
\end{figure}


Layer swap is the design closest to riders: both grow the batch, and both are
bounded by the KV the link can move within a step. They differ in which
directions of the link they use. A swapped layer is written back after
use~\cite{pie}, so layer swap moves as many bytes from device to host as from
host to device, and its swap-out runs in the slower direction, 371 rather
than 419\,GB/s on GH200~\cite{pie}. A rider returns only its new token. The
link is rarely the decode step's alone: saving KV for later reuse,
preemption and co-located work all write to the host, and Pie itself notes
that host--device traffic ``can occur unpredictably due to concurrent
tasks''~\cite{pie}.

We simulate decode in steady state: requests arrive already prefilled, decode
and leave, and every design completes the same requests, admitted in arrival
order. Step times come from the capacity model with the overheads measured for
\sysname{}. Pie is modelled from its own description~\cite{pie}: the
write-back, the slower swap-out and a swap buffer of two layers in HBM. It is
given an ideal split and \sysname{}'s overheads, since its paper reports none.
Other traffic then takes a share of one link direction (\cref{fig:linkload}).
With no other traffic, riders lead layer swap by a median of 3.4\% over 36
workloads. As traffic to the host grows, riders keep their gain, and layer
swap loses most of its own: at half the direction taken, riders lead in all 36
workloads, by a median of 11.7\% and up to 20\%. Traffic toward the device
cuts the other way. Riders depend on that direction alone, and at the same
load they trail layer swap by a median of 1.3\%.
% src: paper/model/steady_state.py (bg_d2h, bg_h2d), steady_bg.py ->
%      model/out/steady_bg.jsonl, overheads charged. hetero/pie: none med +3.4
%      (25 of 36 ahead, -9.3..+8.5); d2h 0.5 med +11.7, 36/36, +0.8..+19.6;
%      h2d 0.5 med -1.3, 12/36, -16.8..+2.7. Pie: s3.3 ("another two layer
%      computation times to swap out"; "both CPU and GPU use the memory space
%      of two layers as a swapping buffer"); s2.4 ("419 GB/s from CPU to GPU
%      and 371 GB/s from GPU to CPU"); s3.4.1 (quoted); s5.5 controller within
%      2% of the best, so the oracle split is generous to Pie. Only the rate
%      lost to other traffic is modelled, not the stall a late swap-out causes
%      through the two-layer buffer. Finished-request KV saved for reuse alone
%      is < 1% of the link at the target (n*S*b/O per step); the larger shares
%      need co-located work.

